% =====================================================================
%  표 2 — 최종 시스템 성능 평가 (사용자 논문에서 2번째 표)
%  컴파일: tectonic -X compile table2_performance.tex
%  본문 삽입: \begin{center}...\end{center} 안의 \captionof + tabular 를 쓰거나,
%  table 환경으로 감싸 쓰면 됨. (논문 자체 번호매김 사용 시 setcounter 불필요)
% =====================================================================
\documentclass[a4paper,11pt]{article}
\usepackage[margin=2.5cm]{geometry}
\usepackage{booktabs}
\usepackage{array}
\usepackage{kotex}
\usepackage{caption}
\captionsetup{labelfont=bf, font=small, skip=8pt}
\renewcommand{\tablename}{표}
\captionsetup{labelsep=period}
\setcounter{table}{1}
\pagestyle{empty}

\begin{document}
\null\vfill
\begin{center}
\captionof{table}{최종 OpenVLA 기반 폐사체 수거 시스템의 성능 평가 결과($n=50$).
개발된 시스템은 50회 반복 평가에서 작업 성공률 90.0\%, 파지 성공률 92.0\%를
달성하였으며, 성공한 모든 사례에서 폐사체를 바구니 중심 5\,cm 이내에 투하하였다.}
\label{tab:performance}

\renewcommand{\arraystretch}{1.3}
\begin{tabular}{l c}
\toprule
평가 항목 & 결과 \\
\midrule
Task 성공률            & 45/50~~(90.0\%) \\
Grasp 성공률           & 46/50~~(92.0\%) \\
투하 정밀도(중앙값)     & 1.34\,cm \\
투하 정밀도(평균)       & 1.47\,cm \\
최대 오차              & 3.93\,cm \\
바구니 중심 $\le 5$\,cm & 45/45~~(100\%) \\
\bottomrule
\end{tabular}
\end{center}
\vfill\null
\end{document}
